\section{Background and Related Work}
\label{sec:related}

The idea that a machine should read code for defects is old, and its
history is a sequence of solutions each of which created the next
problem. The earliest security-relevant static analysers were syntactic:
lint-style checks and, later, bug-pattern detectors such as
FindBugs~\cite{ayewah2008}. They were cheap to run and caught real
mistakes, which established that automated inspection was worthwhile at
all. Their weakness was equally clear, a purely local pattern cannot see
how a tainted value flows across functions, so the field moved to
inter-procedural data-flow and taint analysis. Coverity showed that this
could be done over hundreds of millions of lines of industrial
code~\cite{bessey2010}, Infer brought a compositional formulation to a
comparable scale~\cite{distefano2019}, and Chess and McGraw framed the
whole activity as a security discipline rather than a quality
one~\cite{chess2004}. Depth, however, came with two costs: substantial
engineering to deploy, and a false-positive burden large enough that
managing it became a product in itself. Studies of practitioners found
that developers abandon static analysers not because the tools lack
power, but because results arrive out of context, the noise is high, and
triage does not fit the way people work~\cite{johnson2013,nachtigall2022}.

The response to the rule-authoring bottleneck was to make rules a
first-class, shareable artefact. Yamaguchi \textit{et~al.}\ modelled a
program as a code property graph so that a vulnerability could be
expressed as a graph query~\cite{yamaguchi2014}, and CodeQL turned such
queries into a declarative language with a community
library~\cite{avgustinov2016}; lighter open tools such as
Semgrep~\cite{semgrep} and Bandit~\cite{bandit} put custom rules in reach
of ordinary developers. This solved shareability but not the underlying
limit: a rule, however elegant, encodes a pattern the author already
anticipated, and every new class of weakness needs a new rule. In
parallel, a different line of work summarised a program's health as a
number rather than a list of defects. McCabe's cyclomatic
complexity~\cite{mccabe1976} correlates with defect density and test
effort, and the Maintainability Index folded it together with size and
volume into a single figure~\cite{coleman1994}. These metrics are useful
as a triage signal, but a scalar does not localise or explain a
vulnerability, it only says where to look harder.

Because static reasoning is necessarily incomplete, execution-based
methods developed alongside it. Fuzzing became the dominant automated
bug-finding technique for native code~\cite{manes2021}, and container
isolation made it routine to run unknown or untrusted software with
bounded privileges~\cite{merkel2014docker}. Dynamic analysis finds bugs
that static analysis cannot, but it needs harnesses, it needs time, and
it covers business logic poorly, so it complements rather than replaces
the static passes.

The next shift was semantic. LLMs trained on source code could, for the
first time, reason about a program's intent rather than its
syntax~\cite{chen2021codex}, and researchers quickly turned them toward
security. Pearce \textit{et~al.}\ measured how often code \emph{generated}
by an assistant is insecure~\cite{pearce2022} and then showed that an LLM
can \emph{repair} vulnerabilities zero-shot~\cite{pearce2023}; Khoury
\textit{et~al.}\ found that a large fraction of ChatGPT-produced snippets
contain exploitable flaws~\cite{khoury2023}. Learned detectors such as
Devign~\cite{zhou2019devign} and LineVul~\cite{fu2022linevul} predicted
vulnerable functions and lines directly. The capability was real, but two
new problems came with it. First, the reported accuracy of learned
detectors proved highly sensitive to the quality of the training and
evaluation data, and drops sharply across
projects~\cite{steenhoek2023,croft2023}, which argues for keeping a human
in the loop rather than trusting a verdict. Second, and more
fundamental for practice, the models capable of this reasoning are
reached through hosted APIs, so using them re-introduces exactly the
confidentiality and cost problems that on-premises static analysis had
avoided.

The analysis target then moved as well. As applications began embedding
LLMs, prompt injection emerged as a weakness class in its own
right, direct manipulation of the instruction stream~\cite{perez2022ignore}
and indirect injection through retrieved content~\cite{greshake2023}, and
was codified in the OWASP Top~10 for LLM
Applications~\cite{owaspllm2025}. Retrieval-augmented
generation~\cite{lewis2020rag} is especially exposed, because it adds a
document loader, an embedding store, and a reranker, each a fresh sink for
untrusted input. An analyser built today has to cover this class, not
just the classical one.

The final piece is what makes a different design possible. Open model
families such as LLaMA and Llama~2~\cite{touvron2023llama,touvron2023llama2},
together with efficient quantised inference runtimes like
Ollama~\cite{ollama}, now let a multi-billion-parameter model run on a
laptop at interactive speed. The technical reason the semantic step had
to live in the cloud has therefore gone away.

\cs\ is the point where these threads meet. Prior work advanced one axis
at a time and paid for it on another: depth of analysis at the expense of
usability, semantic power at the expense of confidentiality and cost,
breadth of coverage at the expense of integration. \cs\ instead runs
pattern-based SAST, a cyclomatic-complexity engine, an on-device
\mbox{Llama~3.2} semantic reviewer, and a containerised dynamic pass
together on the developer's machine, behind one interface and one
transparent \phindex, with a hard no-egress boundary. It gives up
frontier-model scale to do so, and manages that trade by keeping the
model advisory and always shown next to a corroborating deterministic
rule and the offending source. \Cref{tab:litreview} summarises the works
above, the advance each made, and the limitation \cs\ is designed to
carry forward.

\begin{table*}[!tbp]
  \centering
  \caption{The evolution of automated code security analysis: what each
           line of work advanced, and the limitation it left open.}
  \label{tab:litreview}
  \renewcommand{\arraystretch}{1.25}
  \small
  \begin{tabularx}{\textwidth}{@{}p{0.20\textwidth} X X@{}}
    \toprule
    \textbf{Study / system} & \textbf{Advance and key findings} &
    \textbf{Limitation left open} \\
    \midrule

    FindBugs; bug patterns \cite{ayewah2008} &
      Simple syntactic patterns find real defects at very low cost;
      established automated inspection as worthwhile. &
      Local patterns cannot follow a tainted value across
      procedures; shallow coverage of semantic bugs. \\

    Coverity at scale \cite{bessey2010} &
      Inter-procedural analysis runs on $10^8$-line industrial
      codebases and finds serious defects in the field. &
      Heavy deployment effort; false-positive triage becomes a
      product in itself. \\

    Infer / compositional analysis \cite{distefano2019} &
      Compositional summaries make deep analysis scale and fit CI. &
      Targets specific bug classes; still rule-/spec-bound. \\

    Developer studies \cite{johnson2013,nachtigall2022} &
      Identify \emph{why} SAST is abandoned: out-of-context results,
      noise, poor workflow fit. &
      Diagnose the adoption problem but propose no tool. \\

    Code property graphs \cite{yamaguchi2014} &
      Represent code as a graph so a vulnerability is a query;
      reusable across a codebase. &
      Requires expert query authoring for each weakness class. \\

    CodeQL \cite{avgustinov2016} &
      Declarative, shareable security queries with a community
      library. &
      A rule encodes a pattern the author anticipated; cannot infer
      intent. \\

    Cyclomatic complexity;
    Maintainability Index \cite{mccabe1976,coleman1994} &
      A structural-health scalar that correlates with defect density
      and test effort. &
      A number does not localise or explain a vulnerability. \\

    Fuzzing survey \cite{manes2021} &
      Execution finds bugs static analysis misses; now the dominant
      automated technique for native code. &
      Needs harnesses and time; weak coverage of business logic. \\

    Codex / LLMs for code \cite{chen2021codex} &
      Models trained on code reason about \emph{intent}, not just
      syntax. &
      General-purpose, not security-tuned; capable variants are
      hosted only. \\

    Copilot / ChatGPT code security
    \cite{pearce2022,khoury2023} &
      Quantify how often assistant-\emph{generated} code is
      insecure. &
      Evaluate the generation side, not a defensive analysis tool;
      prompts still leave the machine. \\

    Devign; LineVul
    \cite{zhou2019devign,fu2022linevul} &
      Learned models predict vulnerable functions / lines directly. &
      Accuracy tied to noisy datasets; weak cross-project
      generalisation. \\

    Zero-shot vulnerability repair \cite{pearce2023} &
      A capable LLM can repair vulnerabilities with no
      task-specific training. &
      Needs a frontier (hosted) model; unreliable without an oracle. \\

    Data-quality critiques
    \cite{steenhoek2023,croft2023} &
      Show DL vulnerability-detection results are inflated by
      label/duplication issues. &
      Argue against trusting a verdict; motivate human-in-the-loop. \\

    Prompt injection; OWASP LLM Top~10
    \cite{perez2022ignore,greshake2023,owaspllm2025} &
      Define a new weakness class (direct and retrieval-borne) that
      analysers must now cover. &
      RAG systems \cite{lewis2020rag} are especially exposed; the
      analysis target itself has moved. \\

    Open quantised models; Ollama
    \cite{touvron2023llama,touvron2023llama2,ollama} &
      Multi-billion-parameter models run on a laptop at interactive
      speed. &
      Smaller than frontier models; recall of subtle,
      cross-function bugs is lower. \\

    \midrule
    \textbf{\cs\ (this work)} &
      Runs pattern SAST, complexity metrics, an on-device
      \mbox{Llama~3.2} reviewer, and a container sandbox together
      behind one interface and one transparent score, with a hard
      no-egress boundary. &
      Trades frontier-model scale for locality; the model is kept
      advisory and corroborated by deterministic rules. \\
    \bottomrule
  \end{tabularx}
\end{table*}

